\subsubsection{The Arzela-Ascoli Theorem}\label{subsec:arzela-ascoli-theorem}

The Arzela-Ascoli theorem states that sequences of functions on a compact set which satisfy certain conditions have uniformly convergent subsequences. Results of this nature are often useful in existence proofs; the thing which must be proved to exist is the limit of the convergent subsequence. To state the theorem, we need the following definition.

\begin{Definition}
\label{def:equicontinuity}
Let $f_m$ be a sequence of real-valued functions defined in a subset $D$ of $\mathbb{R}^n$. Let $\bm{x}\in D$. The sequence is called \textbf{equicontinuous} at $\bm{x}$ if, for every $\epsilon>0$, there is a $\delta>0$, independent of $m$, such that $|f_m(\bm{y})-f_m(\bm{x})|<\epsilon$ for $\bm{y}\in D$ with $|\bm{y}-\bm{x}|<\delta$.
\end{Definition}

If the sequence $f_m$ is equicontinuous at each point of a compact set $S$, it is uniformly equicontinuous, i.e., $\delta$ in the definition above can be chosen independently of $\bm{x}\in S$ (cf. Problem~\ref{ex:compact-equicontinuity-uniform}; it is not necessary that $D=S$). We note that a sequence of functions is equicontinuous at $\bm{x}$ if there exists a bound (independent of $m$) for the derivatives in some neighborhood of $\bm{x}$.

\begin{Theorem}{Arzela-Ascoli}
\label{thm:arzela-ascoli}
Let $f_m$ be a sequence of real-valued functions defined on a compact subset $S$ of $\mathbb{R}^n$. Assume that there is a constant $M$ such that $|f_m(\bm{x})|\leq M$ for every $m\in\mathbb{N}$ and every $\bm{x}\in S$. Moreover, assume that the sequence $f_m$ is equicontinuous at every point of $S$. Then there exists a subsequence which converges uniformly on $S$.
\end{Theorem}

We remark that (with rather obvious modifications) the theorem and proof can be extended to the case where $S$ is a compact set in an abstract topological space and the values of $f_m$ are in a complete metric space.

\begin{Proof}
Let $\bm{x}_i$, $i\in\mathbb{N}$ be a sequence of points that is dense in $S$. The sequence $f_m(\bm{x}_1)$ is bounded; hence it has a convergent subsequence. That is, we can choose a sequence $m_{1j}$ such that $f_{m_{1j}}(\bm{x}_1)$ converges as $j\to\infty$. Similarly, we can choose a subsequence $m_{2j}$ of the sequence $m_{1j}$ such that $f_{m_{2j}}(\bm{x}_2)$ converges. Since $m_{2j}$ is a subsequence of $m_{1j}$, $f_{m_{2j}}(\bm{x}_1)$ converges as well. Next, we choose a subsequence $m_{3j}$ of the sequence $m_{2j}$ such that $f_{m_{3j}}$ converges also at $\bm{x}_3$. We proceed in this manner ad infinitum. Finally, consider the "diagonal" sequence $f_{m_{jj}}$. Except for the first $i-1$ terms, $m_{jj}$ is a subsequence of $m_{ij}$; hence $f_{m_{jj}}(\bm{x}_i)$ converges for every $i\in\mathbb{N}$. To simplify notation, we shall set $g_j=f_{m_{jj}}$ in the following.

To conclude the proof, we show that the sequence $g_m$ is uniformly Cauchy. Let $\epsilon>0$ be given. The $g_m$, being a subsequence of the $f_m$, are uniformly equicontinuous on $S$; hence there is a $\delta>0$ such that $|g_m(\bm{y})-g_m(\bm{x})|<\epsilon/3$ whenever $|\bm{y}-\bm{x}|<\delta$. Since $S$ is compact, there is a $K\in\mathbb{N}$ such that for every $\bm{x}\in S$ there exists $i\in\{1,\ldots,K\}$ with $|\bm{x}_i-\bm{x}|<\delta$. Now choose $N$ large enough so that $|g_m(\bm{x}_i)-g_k(\bm{x}_i)|<\epsilon/3$ for $m,k>N$ and every $i\in\{1,\ldots,K\}$. For $m,k>N$ and arbitrary $\bm{x}\in S$, we now have
\begin{equation}\label{eq:arzela-ascoli-uniform-cauchy-estimate}
\begin{aligned}
|g_m(\bm{x})-g_k(\bm{x})|
&\leq|g_m(\bm{x})-g_m(\bm{x}_i)|\\
&\quad+|g_m(\bm{x}_i)-g_k(\bm{x}_i)|+|g_k(\bm{x}_i)-g_k(\bm{x})|\\
&<\epsilon,
\end{aligned}
\end{equation}
for some $i\in\{1,\ldots,K\}$.
\end{Proof}

Below, we shall have to apply the Arzela-Ascoli theorem to sequences of harmonic functions. In this particular case, we have the following result.

\begin{Theorem}
\label{thm:harmonic-sequence-compactness}
Let $\Omega$ be a domain in $\mathbb{R}^n$. Let $f_m$ be a sequence of harmonic functions on $\Omega$ which is uniformly bounded, i.e., $|f_m(\bm{x})|\leq M$ for every $\bm{x}\in\Omega$ and $m\in\mathbb{N}$. Then $f_m$ has a subsequence which converges to a harmonic function on $\Omega$, uniformly on compact subsets of $\Omega$.
\end{Theorem}

\begin{Proof}
Let $\Omega_k=\{\bm{x}\in\Omega\mid|\bm{x}|\leq k,\ \operatorname{dist}(\bm{x},\partial\Omega)\geq1/k\}$. Then $\Omega_k$ is compact and $\Omega=\bigcup_{k=1}^\infty\Omega_k$. If $\bm{x}$ is in $\Omega_k$, then, in a neighborhood of $\bm{x}$, we can represent $f_m$ by Poisson\textquotesingle s formula, using a ball centered at $\bm{x}$ with radius $1/2k$. This yields uniform estimates for the derivatives of $f_m$ on $\Omega_k$. In particular, the $f_m$ are equicontinuous on $\Omega_k$. By the Arzela-Ascoli theorem, we can extract a uniformly convergent subsequence on each $\Omega_k$, and using a diagonal argument similar to that used in the proof above, we obtain a subsequence converging on all of $\Omega$, uniformly on each $\Omega_k$.

It remains to show that the limit of this subsequence is harmonic. To see this, we simply note that harmonic functions are characterized by the property that they obey Poisson\textquotesingle s formula on every ball. If Poisson\textquotesingle s formula is obeyed by all functions in a uniformly convergent sequence, it also holds for the limit.
\end{Proof}
